\chapter{Discussion}
\label{ch:discussion}

\section{What the geometry decides}

The most consequential finding of this project is not a control result but a
geometric one: the two arms, as mounted and as parked, share no reachable
workspace. That single fact removes an entire class of task from the
programme, and it does so in a way no amount of planning, control or autonomy
can address, because the constraint is on where the arms can be at all rather
than on how they get there.

It is worth being precise about what follows from it and what does not. Tasks
requiring the two arms to meet at a point are blocked. Tasks requiring two
distinct points a fixed distance apart are not, and were shown to be feasible
in a band at chest height. The distinction was initially missed, and a
bimanual programme was briefly declared blocked on the strength of a test that
answered the stronger question. The correction matters more than the original
error: the requirement a task actually imposes should be tested, rather than a
convenient proxy for it.

\section{Clearance and reach are in conflict, and the conflict is structural}

The clearance result in Chapter~\ref{ch:results} deserves emphasis because it
is easy to read as a tuning problem. It is not. The planner's collision check
and the deployed clearance floor ask different questions --- contact versus
margin --- and a scenario can satisfy the first while being refused by the
second. Every chest-band scenario in this project is in exactly that position.

Attempting to resolve it by moving the working volume forward fails, and the
failure is instructive: clearance improves and reachability collapses, because
the volume moves towards the edge of the arm's reach as it moves away from the
wearer's body. The axis along which the conflict can be resolved is lateral,
not fore-and-aft. That was not obvious in advance and was found only because
the probe measured which link was closest, rather than only how close it was.

The ethical position is one-directional. The floor exists to protect a person
who did not choose the motion, so the task moves and the floor does not. That
constraint is recorded rather than negotiated.

\section{Instrumentation failed more often than the system}

A recurring pattern across this project, and one which shaped its methodology
more than any design decision, is that surprising measurements were more often
faults in the measuring instrument than in the system being measured.

The instances are numerous enough to be a finding rather than an anecdote. A
noise floor of zero on every sensor channel was the recorder oversampling a
slower sensor. A detection rate near zero was a renderer producing images
outside the detector's training distribution. A workspace statistic that
varied by a factor of two between identical runs was a summary statistic
sensitive to a rare tail event. A set of verification clips reported as
failing were failing a colour test calibrated against the requested marker
colour rather than the rendered one. A carried object reported as stationary
was a centroid dragged towards fixed overlay text. In each case the system was
behaving correctly and the instrument was not.

The methodological response --- requiring every analysis routine to recover a
known answer before it is trusted with real data --- is, in the author's view,
the most transferable contribution of the work. It is also incomplete: the
rule was adopted partway through and the earlier measurements have not all
been re-derived under it.

\section{Silent failure as the dominant failure mode}

The safety work in this project consistently found that the problem was not
that mechanisms failed to act, but that acting and not acting looked identical
from outside. The dead-man defect is the clearest example: the mechanism was
correct, its trigger condition was correct, and it could not fire on the
failure that actually occurs, because the upstream node degraded by
republishing stale data rather than by falling silent. Every test of it had
passed, because every test removed the cable.

The general lesson is that a test which exercises a failure mode different
from the one that occurs provides confidence without providing safety. The
system now distinguishes data that is absent from data that is stale, and
treats a guard that has stopped reporting as being in an unknown state rather
than a clear one.

\section{Why no end-to-end learned policy}

A vision--language--action model was considered and deliberately not used, for
three reasons that are worth recording because the decision would otherwise
look like conservatism.

First, embodiment: available models are trained on datasets that contain no
backpack-mounted supernumerary configuration, and there is no demonstration
data for this platform with which to fine-tune. Second, and decisively, such a
model emits actions directly, whereas every safety mechanism in this system
sits between a commanded \emph{pose} and the arm. A policy emitting joint
targets bypasses collision-aware inverse kinematics, the clearance floor and
the step guard. On a robot mounted beside a person's head that is the wrong
trade. Third, a failure in the modular pipeline is attributable to a named
stage with its inputs logged, whereas a failure of a learned policy is a
number that was wrong.

\section{Limitations}

\subsection{Nothing has been validated on hardware}

The most serious limitation is stated plainly: no result in this thesis was
obtained from a physical arm, a physical master or a human participant. The
simulation uses a mock hardware interface that echoes commands with no
dynamics, which means it cannot answer questions about payload, compliance,
latency under load, or any effect of real actuator behaviour. Where the
simulation has been asked such a question it has returned a null result that
is a property of the mock rather than evidence about the robot.

\subsection{The input is degraded and the baseline depends on it}

Seven of fourteen master-arm channels are unreliable. Every direct
teleoperation condition in the experimental design depends on those channels,
so the baseline against which autonomy would be compared is currently
compromised. Repairing two specific channels was identified as the
highest-value intervention.

\subsection{The experimental programme is designed but unexecuted}

The protocols, measures, counterbalancing and analysis routines exist and have
been exercised with scripted input, but no participant has taken part. The
hypotheses are pre-registered intent. \todo{CONFIRM: whether ethics approval
was sought, and its status, since this determines how the design is described}

\subsection{The commit history is not a development record}

The project's version control history consists of one large initial import and
a single day of subsequent commits, so it cannot support claims about the
order or timing of development. The dated engineering log is the usable
chronology, and it is a laboratory notebook rather than an independent record.
